\section{Kết luận và hướng tiếp}\label{sec:conclusion}

Theo lời tác giả, ViPharmind được làm với trợ lý AI ở mọi khâu; lịch sử git cho thấy nếp spec, plan rõ:
\NumFeatWithPlanPct{}\% commit tính năng nối được với một plan (cận trên, vì quy tắc nối nghiêng về nối
thừa), và commit tính năng có hay không có plan bị sửa lại gần như ngang nhau (RQ1, RQ2).

Khi đề đã nêu đủ tên và chữ ký, quy trình đặc tả không cho thấy đạt test ẩn nhiều hơn ở việc nhỏ (H1, H5
chưa kết luận được) mà tốn gấp \NumResTokenRatioMin{} đến \NumResTokenRatioMax{} lần token; dấu hiệu có ích duy nhất là Codex trên ba nhiệm vụ lớn nhất, và
\NumResSonnetSpecTestsAnyPct{}--\NumResCodexSpecTestsAnyPct{}\% số lượt theo đặc tả để lại test do tác tử tự viết. \NumLabOutside{} trên \NumLabRuns{} lượt
hỏng của Claude đến từ chi tiết đề không nêu nhưng test ẩn kiểm, hoặc từ việc nối với mã có sẵn, ở cả hai
cách làm như nhau. Claude và Codex lệch hướng nhau nhưng ở tầng nhỏ mọi khoảng tin cậy đều chứa 0 (RQ4); H4 chưa kết luận
được (RQ5).

Với đề chỉ một hai câu, thí nghiệm thăm dò E2 gợi ý rằng đề rõ quan trọng hơn quy trình: bỏ tên và chữ
ký làm điểm danh sách hành vi giảm \NumETwoXScoreShiftSpec{} đến \NumETwoXScoreShiftVibe{} đpt ở cả hai
cách. Khi đề sơ sài, điểm của nhánh đặc tả cao hơn vibe \NumETwoPooledSmallPp{}
đpt ở tầng nhỏ [\NumETwoPooledSmallLowPp{}; \NumETwoPooledSmallHighPp{}]; gộp mọi tầng
là \NumETwoPooledAllPp{} đpt, khoảng tin cậy chứa 0. E1 (Haiku) cho kết quả cùng hướng. \textbf{Điểm chính:} khi đề đã chi tiết, đặc tả
tốn thêm mà không đạt nhiều hơn; khi đề sơ sài, nó có thể bù một phần ở việc nhỏ,
rõ nhất với Sonnet, nhưng không thay được thông tin đề thiếu (thăm dò, cần kiểm lại bằng thí nghiệm đăng
ký trước). Hướng tiếp: nhánh đặc tả có người duyệt, thêm dự án và nhiệm vụ cỡ tính năng.
